\section{Discussion: two failure modes, two repairs, and what remains open}
\label{sec:discussion}

\subsection{The shape of the answer}

The finite picture in Table~\ref{tab:finite} is neither a proof nor a refutation
of a single sentence. It is a map of an interface. Bridging conditions are a
\emph{structural} hypothesis; maximum likelihood is an \emph{optimization}
objective; the published sentence asks whether the first implies the second. Our
study shows that the answer depends on choices the sentence does not make, and
that the mathematically clean content lives in two distinct failure modes.

\paragraph{Structural/model failure.}
Without a repeat/read-length boundary the genome need not be identifiable at all
(Example~\ref{ex:ABAB}). Bridging conditions repair the structural ambiguity, and
Theorem~\ref{thm:rigidity} shows that on the same-length oriented Section~6.2
slice they are strong enough to make the truth's spectrum the only positive
circulation --- an exact rigidity, not just a reconstruction claim. Requiring
candidate length equal to the true length is, however, a strong extra axiom. It
is not supplied by Medvedev and Brudno's known-size parameter, and
Example~\ref{ex:variable-length} shows the conclusion fails as soon as it is
dropped.

\paragraph{Finite-sampling failure.}
Even structurally admissible genomes can lose to the truth's competitors because
empirical read frequencies fluctuate (Example~\ref{ex:frequencies}). This is a
different mechanism and needs a different repair. Population maximum likelihood
removes it at the maximizer level by construction
(Lemma~\ref{lem:gibbs}), and for the oriented primitive P2 class it also yields
uniqueness (Theorem~\ref{thm:population}).

The resulting conceptual progression is:

\begin{quote}
naive ML intuition
$\to$ structural repeat obstruction
$\to$ the literature's repeat/bridging boundary
$\to$ boundary-conditioned finite ML
$\to$ a positive same-length slice but negative variable-length neighbors
$\to$ replacing privileged true-length knowledge by intrinsic candidate checks
$\to$ finite failure surviving as sampling fluctuation
$\to$ population ML removing the fluctuation
$\to$ P2 plus primitivity restoring oriented circular uniqueness.
\end{quote}

This ordering matters: population data is not introduced to patch repeat
ambiguity, and intrinsic candidate checks are not introduced as assumptions of
the historical papers. They repair weaknesses exposed in sequence.

\subsection{What remains open}

\paragraph{Source questions.}
Which Medvedev--Brudno object the 2016 sentence denotes; which read-type
convention (and hence which equivalence) it intends; whether candidate length is
restricted; whether the conclusion is maximizer or uniqueness; and whether the
unretrieved publisher supplement contains a definition that narrows any of
these. These are not bookkeeping: each changes which formal statement the
published sentence asserts.

\paragraph{Mathematical questions.}
The same-length Section~6.2 statement under the per-occurrence strengthening
($d_D\ge x$) remains open in the searched scope. The single-strand reading of the
variable-length Section~6.2 statement remains open (bounded zero evidence). The
population uniqueness theorem is now supported by the P2-specific gcd-one
argument in Lemma~\ref{lem:scaling} together with the external
Bresler--Bresler--Tse uniqueness theorem; this chain remains source-supported
mathematics rather than a kernel-checked Lean theorem. The population model for
\emph{molecule} read types --- as opposed to oriented types --- is not addressed
here.

\paragraph{What would count as settlement.}
A positive settlement of the published question is a Lean proof of a formally
justified version of the published implication after an argument fixes the axes
of Table~\ref{tab:axes}. A negative settlement is a mathematically valid
counterexample satisfying the faithfully formalized bridging hypotheses while
violating the faithfully formalized likelihood conclusion. A computational
discovery is useful but not a settlement until the finite instance is
kernel-checked and the source correspondence is documented.

\subsection{Relation to the literature}

The closest published result is older than the 2016 question:
Bresler, Bresler, and Tse prove a \emph{necessary} condition --- if a
problematic interleaved pair or triple repeat has all copies unbridged, another
same-length genome has the same read likelihood. Shomorony et al.\ strengthen
the hypothesis to sufficient conditions for exact reconstruction and then ask
whether those conditions also imply ML optimality. A 2025 correctness-guarantee
paper still presents the information/correctness line and the
maximum-likelihood line as separate \cite{salmela2025}, which is meaningful
evidence that the 2016 question had not become a standard solved result. Our
negative witnesses are consistent with that picture: the implication is not
merely unproved but false under several precise readings, while the positive
same-length and population results show where the intuition does survive.

\subsection{Conclusion}

The 2016 question is best read not as one theorem awaiting one proof but as an
interface between a bridging hypothesis and a maximum-likelihood formulation
whose candidate, representation, length, and tie choices all matter. We have made
those choices explicit, proved a positive same-length rigidity theorem, refuted
several precise readings by strict kernel-checked witnesses, and presented a
repaired population model that restores oriented circular uniqueness under
primitivity and a complete-spectrum condition. The finite source-faithful
question remains distinct from the repaired result, and we have marked exactly
where the boundary lies.
